\section*{Chapitre \ref{ch:m3:margin-liquidation-and-loss-allocation} --- Marge, liquidation et répartition des pertes}

\begin{solution}{exo:m3:margin-liquidation-and-loss-allocation:1}
Marge de 5 par unité~: liquidation à $95/0.995 = 95.48$, faillite à 95.
\end{solution}

\begin{solution}{exo:m3:margin-liquidation-and-loss-allocation:2}
Non~: en \omterm{def:m3:margin-liquidation-and-loss-allocation:modes}{marge croisée}, la totalité du solde soutient la position, de sorte qu'elle n'est liquidée que lorsque les fonds propres totaux du compte
tombent à leur marge de maintenance, bien plus loin~; mais une baisse suffisamment forte absorbe la totalité du solde.
\end{solution}

\begin{solution}{exo:m3:margin-liquidation-and-loss-allocation:3}
La marge du trader, la vente du \omterm{def:m3:margin-liquidation-and-loss-allocation:engine}{moteur de liquidation} (excédent ou déficit), le \omterm{def:m3:margin-liquidation-and-loss-allocation:fund}{fonds d'assurance}, l'\omterm{def:m3:margin-liquidation-and-loss-allocation:adl}{auto-deleveraging} (ADL) des
positions opposées bénéficiaires et, sur certaines plateformes ou dans le passé, une \omterm{def:m3:margin-liquidation-and-loss-allocation:adl}{perte mutualisée}.
\end{solution}

\begin{solution}{exo:m3:margin-liquidation-and-loss-allocation:4}
Linéaire 80.40, inverse 83.75. La marge de la position longue inverse est libellée en coin, qui perd de sa valeur à mesure que le cours baisse, de sorte que les fonds propres
diminuent plus vite que la seule perte en dollars de la position.
\end{solution}

\begin{solution}{exo:m3:margin-liquidation-and-loss-allocation:5}
$1\%/(1 - 0.8) = 5\%$.
\end{solution}

\begin{solution}{exo:m3:margin-liquidation-and-loss-allocation:6}
En cas de krach, sa position short sur \omterm{def:m3:perpetual-futures:perp}{contrat perpétuel} compte parmi les plus profitables et peut être clôturée par \omterm{def:m3:margin-liquidation-and-loss-allocation:adl}{auto-deleveraging}, ce qui laisse
sa position longue au comptant non couverte au moment où les cours chutent le plus vite. Elle surveille l'indicateur ADL, maintient un faible effet de levier sur la position short
(ce qui abaisse son rang), répartit la couverture entre plusieurs plateformes et prévoit de se recouvrir immédiatement.
\end{solution}

\begin{solution}{exo:m3:margin-liquidation-and-loss-allocation:7}
1.65~\% avec un effet de levier plafonné à 20 fois et 5.37~\% à 10 fois, contre 1.11~\% sans plafond~: la concentration de \omterm{def:m3:margin-liquidation-and-loss-allocation:prices}{prix de liquidation} à fort
effet de levier près du prix d'entrée est ce qui rend un petit choc dangereux~; les plafonds d'effet de levier protègent donc le marché, et pas seulement
les traders.
\end{solution}

\begin{solution}{exo:m3:margin-liquidation-and-loss-allocation:8}
Avec un impact permanent linéaire, le prix final est le même quelle que soit la taille du lot
(voir la \cref{prop:m3:margin-liquidation-and-loss-allocation:kappa}), et il en va de même du total du \omterm{def:m3:margin-liquidation-and-loss-allocation:fund}{fonds d'assurance}, puisque le prix d'exécution moyen pondéré par les unités
est le même~; le tutoriel vérifie les deux. La réponse ne change que si l'impact est non linéaire ou si le carnet d'ordres se reconstitue entre les lots, et alors un pas plus fin
relève d'un autre modèle, non d'un modèle plus précis.
\end{solution}

\begin{solution}{pb:m3:margin-liquidation-and-loss-allocation:1}
\textbf{1.} Marge de 1\,000~; liquidation à $24\,000/(250 \times 0.995) = 96.48$~; faillite à 96. \textbf{2.} $250
\times 0.5 = 125$. \textbf{3.} Il reçoit $250 \times 0.8 = 200$. \textbf{4.} Vendue au-dessus du \omterm{def:m3:margin-liquidation-and-loss-allocation:prices}{prix de faillite}, ce qui reste
de la marge du trader revient au fonds, ce qui explique la croissance des fonds en temps normal. \textbf{5.} La dernière transaction peut être déplacée
à bon compte pendant un instant~; le \omterm{def:m3:perpetual-futures:index}{prix de marquage}, construit sur un indice multiplateforme, ne le peut pas. \textbf{6.} 1.11~\%. \textbf{7.} Une~: chaque point de
baisse à ce niveau déclenche des liquidations qui provoquent une nouvelle baisse supérieure à un point, de sorte que la cascade traverse toute la zone.
\textbf{8.} 5.37~\% avec le plafond~; avec le carnet d'ordres plus profond, le cours n'atteint une baisse de 10~\% qu'après un choc de 5.55~\%, et sans saut, puisque $\kappa$ est alors inférieur à un partout. \textbf{9.} Avec un impact permanent linéaire, le prix ne dépend que des unités vendues, et non
de leur ordre. \textbf{10.} Un impact non linéaire, un carnet d'ordres qui se reconstitue entre les lots, ou des délais des moteurs qui laissent les prix sauter
entre les lots. \textbf{11.} Autour de 83.83, avec 3.21~millions de dollars américains de déficits. \textbf{12.} Environ 1.21~millions de dollars américains~: les
déficits moins les 2~millions de dollars américains du fonds (et quelques milliers de dollars d'excédents). \textbf{13.} Les positions short affichant le pourcentage de
P\&L le plus élevé multiplié par l'effet de levier effectif~: les gagnants les plus profitables et les plus endettés. \textbf{14.} Sa position short est clôturée au
\omterm{def:m3:margin-liquidation-and-loss-allocation:prices}{prix de faillite}~; elle conserve donc le P\&L jusqu'à ce prix, mais perd la couverture et le reste du mouvement, et doit se recouvrir
sur un marché en baisse. \textbf{15.} Elle répartirait la perte entre tous les gagnants au prorata de leurs gains ou de leur risque, de sorte
qu'aucun gagnant ne perde la totalité de sa position. \textbf{16.} Les plafonds relèvent le choc qu'un marché peut absorber (de 1.11~\% à 5.37~\% ici), au
prix de la redirection de la demande à effet de levier vers d'autres plateformes~; une plateforme qui plafonne seule peut perdre le flux sans rendre le
marché plus sûr. \textbf{17.} Non~: un écart de 5~\% à lui seul le consomme. Un fonds doit être dimensionné selon des scénarios de stress portant sur les positions ouvertes,
l'effet de levier et la profondeur, et non l'historique. \textbf{18.} Comme donnée de risque en temps réel~: réduire l'effet de levier ou la taille des positions situées haut dans la
file, et traiter un rang élevé comme le signal qu'une couverture risque de disparaître. \textbf{19.} \emph{Résultat clé~:} un choc critique de
1.11~\%~; après un écart de 5~\%, 3.21~millions de dollars américains de déficits, dont 2~millions de dollars américains provenant du \omterm{def:m3:margin-liquidation-and-loss-allocation:fund}{fonds d'assurance} et environ 1.21~millions de dollars américains
par \omterm{def:m3:margin-liquidation-and-loss-allocation:adl}{auto-deleveraging}. \textbf{20.} La marge du perdant, le \omterm{def:m3:margin-liquidation-and-loss-allocation:fund}{fonds d'assurance} de la plateforme, puis les autres gagnants.
\end{solution}

\begin{solution}{iq:m3:margin-liquidation-and-loss-allocation:1}
Marge $M = qE/10$. Fonds propres $M + q(P - E)$ rejoignent $m\,qP$ en $P = (qE - M)/(q(1 - m)) = 0.9E/(1 - m)$~: environ 9.5~\% sous le
prix d'entrée pour $m = 0.5\%$.
\iqlookfor{fonds propres égaux à la marge de maintenance au \omterm{def:m3:perpetual-futures:index}{prix de marquage}.}
\end{solution}

\begin{solution}{iq:m3:margin-liquidation-and-loss-allocation:2}
Lorsque le \omterm{def:m3:margin-liquidation-and-loss-allocation:fund}{fonds d'assurance} ne peut pas absorber une position en faillite, les positions opposées bénéficiaires sont clôturées au \omterm{def:m3:margin-liquidation-and-loss-allocation:prices}{prix de faillite} ou au
\omterm{def:m3:perpetual-futures:index}{prix de marquage}. Un hedger non endetté peut lui aussi subir l'\omterm{def:m3:margin-liquidation-and-loss-allocation:adl}{auto-deleveraging}~: sa couverture disparaît lors d'un krach, et ses gains sont plafonnés à un prix
qu'il n'a pas choisi.
\iqlookfor{les gagnants supportent la queue~; les couvertures peuvent disparaître.}
\end{solution}

\begin{solution}{iq:m3:margin-liquidation-and-loss-allocation:3}
La densité des \omterm{def:m3:margin-liquidation-and-loss-allocation:prices}{prix de liquidation} près du cours actuel (effet de levier et distribution des prix d'entrée), l'impact de prix des ventes forcées
(profondeur, et reconstitution ou non du carnet), la vitesse des moteurs par rapport à l'arbitrage et aux nouvelles liquidités, et les positions corrélées
entre plateformes~: le multiplicateur local $\kappa$.
\iqlookfor{densité multipliée par l'impact, et franchissement ou non de un.}
\end{solution}

\begin{solution}{iq:m3:margin-liquidation-and-loss-allocation:4}
Liquider tôt et partiellement (réduire vers les paliers inférieurs avant la liquidation complète), fractionner les ordres selon la profondeur du carnet, faire appel à des fournisseurs de liquidité de secours
pour les grandes positions, déclencher sur un \omterm{def:m3:perpetual-futures:index}{prix de marquage} robuste, dimensionner les marges de maintenance selon la taille de la position et la liquidité
de l'actif, et garder le fonds dimensionné selon des scénarios de stress.
\iqlookfor{liquidité partielle, patiente et préalablement convenue.}
\end{solution}

\begin{solution}{iq:m3:margin-liquidation-and-loss-allocation:5}
La \omterm{def:m3:margin-liquidation-and-loss-allocation:modes}{marge croisée} mutualise le collatéral, de sorte que les gains sur certaines cotations compensent les pertes sur d'autres et que moins de capital reste inemployé, mais une
mauvaise position peut emporter tout le compte. La \omterm{def:m3:margin-liquidation-and-loss-allocation:modes}{marge isolée} contient chaque perte, mais exige du capital sur chaque position et liquide
des positions qu'une vision de portefeuille conserverait. Un teneur de marché utilise en général une \omterm{def:m3:margin-liquidation-and-loss-allocation:modes}{marge croisée} ou une marge en portefeuille, assortie de ses propres limites par instrument.
\iqlookfor{efficacité contre confinement, avec des limites internes.}
\end{solution}

\begin{solution}{iq:m3:margin-liquidation-and-loss-allocation:6}
À partir des variations de positions ouvertes par niveau de prix (une hausse des positions ouvertes à un prix suggère des entrées à cet endroit), des \omterm{def:m3:perpetual-futures:perp}{taux de financement} et des ratios
longs--shorts, des distributions d'effet de levier publiées par les plateformes ou estimées à partir des données de marge, et des flux de liquidations a posteriori~; on les combine
en une distribution des entrées multipliée par une distribution d'effet de levier, puis calibrer le modèle sur les cascades observées.
\iqlookfor{entrées déduites des positions ouvertes, effet de levier issu de tout substitut disponible, et calibration.}
\end{solution}
